\QHeading{第1题　口算。}
\par\noindent


\begin{oralgrid}
\OralItem{1}{0.28+0.02}{0.3}
\OralItem{2}{6.5+0.3}{6.8}
\OralItem{3}{3.85-3}{0.85}
\OralItem{4}{5.4+3.6}{9}
\OralItem{5}{3.5-1.5}{2}
\OralItem{6}{2.5+0.5}{3}
\OralItem{7}{1.65-1.35}{0.3}
\OralItem{8}{11+4.72}{15.72}
\OralItem{9}{3.8+1.4}{5.2}
\OralItem{10}{5.7-2.7}{3}
\OralItem{11}{4.3-1.4}{2.9}
\OralItem{12}{10+5.05}{15.05}\end{oralgrid}
\par

\QHeading{口算思路补充（对应上表题号）}

\Answer{上面列出得数，下面按小题给出不同计算思路。计数单位法适用于小数，整数题可比较凑整、拆分或商不变等方法。}
\MethodTitle{第1小题：$0.28+0.02$}

\MethodBlock{口算思路一：按相同计数单位计算}{先看成若干个百分之一，按整数加减，最后还原成小数。}{\begin{aligned}[t]&\phantom{={}}0.28+0.02\\
&=(28+2)\div 100\\
&=30\div 100\\
&=0.3\end{aligned}}
\MethodTitle{第2小题：$6.5+0.3$}

\MethodBlock{口算思路一：按相同计数单位计算}{先看成若干个十分之一，按整数加减，最后还原成小数。}{\begin{aligned}[t]&\phantom{={}}6.5+0.3\\
&=(65+3)\div 10\\
&=68\div 10\\
&=6.8\end{aligned}}
\MethodTitle{第3小题：$3.85-3$}

\MethodBlock{口算思路一：按相同计数单位计算}{先看成若干个百分之一，按整数加减，最后还原成小数。}{\begin{aligned}[t]&\phantom{={}}3.85-3\\
&=(385-300)\div 100\\
&=85\div 100\\
&=0.85\end{aligned}}
\MethodTitle{第4小题：$5.4+3.6$}

\MethodBlock{口算思路一：按相同计数单位计算}{先看成若干个十分之一，按整数加减，最后还原成小数。}{\begin{aligned}[t]&\phantom{={}}5.4+3.6\\
&=(54+36)\div 10\\
&=90\div 10\\
&=9\end{aligned}}
\MethodTitle{第5小题：$3.5-1.5$}

\MethodBlock{口算思路一：按相同计数单位计算}{先看成若干个十分之一，按整数加减，最后还原成小数。}{\begin{aligned}[t]&\phantom{={}}3.5-1.5\\
&=(35-15)\div 10\\
&=20\div 10\\
&=2\end{aligned}}

\MethodBlock{口算思路二：补差法}{从减数先补到一个整数，再补到被减数，两段差相加。}{\begin{aligned}[t]&\phantom{={}}3.5-1.5\\
&=(2-1.5)+(3.5-2)\\
&=0.5+1.5\\
&=2\end{aligned}}
\MethodTitle{第6小题：$2.5+0.5$}

\MethodBlock{口算思路一：按相同计数单位计算}{先看成若干个十分之一，按整数加减，最后还原成小数。}{\begin{aligned}[t]&\phantom{={}}2.5+0.5\\
&=(25+5)\div 10\\
&=30\div 10\\
&=3\end{aligned}}
\MethodTitle{第7小题：$1.65-1.35$}

\MethodBlock{口算思路一：按相同计数单位计算}{先看成若干个百分之一，按整数加减，最后还原成小数。}{\begin{aligned}[t]&\phantom{={}}1.65-1.35\\
&=(165-135)\div 100\\
&=30\div 100\\
&=0.3\end{aligned}}
\MethodTitle{第8小题：$11+4.72$}

\MethodBlock{口算思路一：按相同计数单位计算}{先看成若干个百分之一，按整数加减，最后还原成小数。}{\begin{aligned}[t]&\phantom{={}}11+4.72\\
&=(1100+472)\div 100\\
&=1572\div 100\\
&=15.72\end{aligned}}

\MethodBlock{口算思路二：凑整再减去多加的数}{把第二个数凑成整数，再减去多加的部分。}{\begin{aligned}[t]&\phantom{={}}11+4.72\\
&=11+5-0.28\\
&=16-0.28\\
&=15.72\end{aligned}}
\MethodTitle{第9小题：$3.8+1.4$}

\MethodBlock{口算思路一：按相同计数单位计算}{先看成若干个十分之一，按整数加减，最后还原成小数。}{\begin{aligned}[t]&\phantom{={}}3.8+1.4\\
&=(38+14)\div 10\\
&=52\div 10\\
&=5.2\end{aligned}}
\MethodTitle{第10小题：$5.7-2.7$}

\MethodBlock{口算思路一：按相同计数单位计算}{先看成若干个十分之一，按整数加减，最后还原成小数。}{\begin{aligned}[t]&\phantom{={}}5.7-2.7\\
&=(57-27)\div 10\\
&=30\div 10\\
&=3\end{aligned}}

\MethodBlock{口算思路二：补差法}{从减数先补到一个整数，再补到被减数，两段差相加。}{\begin{aligned}[t]&\phantom{={}}5.7-2.7\\
&=(3-2.7)+(5.7-3)\\
&=0.3+2.7\\
&=3\end{aligned}}
\MethodTitle{第11小题：$4.3-1.4$}

\MethodBlock{口算思路一：按相同计数单位计算}{先看成若干个十分之一，按整数加减，最后还原成小数。}{\begin{aligned}[t]&\phantom{={}}4.3-1.4\\
&=(43-14)\div 10\\
&=29\div 10\\
&=2.9\end{aligned}}

\MethodBlock{口算思路二：补差法}{从减数先补到一个整数，再补到被减数，两段差相加。}{\begin{aligned}[t]&\phantom{={}}4.3-1.4\\
&=(2-1.4)+(4.3-2)\\
&=0.6+2.3\\
&=2.9\end{aligned}}
\MethodTitle{第12小题：$10+5.05$}

\MethodBlock{口算思路一：按相同计数单位计算}{先看成若干个百分之一，按整数加减，最后还原成小数。}{\begin{aligned}[t]&\phantom{={}}10+5.05\\
&=(1000+505)\div 100\\
&=1505\div 100\\
&=15.05\end{aligned}}


\par\noindent

\QHeading{第2题　列竖式计算。}

\Problem{（1）$5.36+2.43=7.79$}
\par


\MethodTitle{方法一：按题意列竖式}
\Answer{\begin{tikzpicture}[x=1em,y=1em,baseline=(current bounding box.center),draw=answerred,text=answerred]
\node[anchor=base,inner sep=0pt] at (0.000,0) {$5$};
\node[anchor=base,inner sep=0pt] at (0.800,0) {$.$};
\node[anchor=base,inner sep=0pt] at (1.100,0) {$3$};
\node[anchor=base,inner sep=0pt] at (1.900,0) {$6$};
\node[anchor=base,inner sep=0pt] at (0.000,-1.1) {$2$};
\node[anchor=base,inner sep=0pt] at (0.800,-1.1) {$.$};
\node[anchor=base,inner sep=0pt] at (1.100,-1.1) {$4$};
\node[anchor=base,inner sep=0pt] at (1.900,-1.1) {$3$};
\node[anchor=base,inner sep=0pt] at (0.000,-2.3) {$7$};
\node[anchor=base,inner sep=0pt] at (0.800,-2.3) {$.$};
\node[anchor=base,inner sep=0pt] at (1.100,-2.3) {$7$};
\node[anchor=base,inner sep=0pt] at (1.900,-2.3) {$9$};
\node[anchor=base,inner sep=0pt] at (-.7,-1.1) {$+$};
\draw[line width=.5pt] (-1,-1.52)--(2.500,-1.52);
\end{tikzpicture}}

\MethodBlock{方法二：拆成整数部分和小数部分}{上方竖式及题目要求的验算仍保留；本法用于核对。}{\begin{aligned}[t]&\phantom{={}}5.36+2.43\\
&=5.36+2+0.43\\
&=7.36+0.43\\
&=7.79\end{aligned}}

\MethodBlock{方法三：按相同计数单位计算}{上方竖式及题目要求的验算仍保留；本法用于核对。先看成若干个百分之一，按整数加减，最后还原成小数。}{\begin{aligned}[t]&\phantom{={}}5.36+2.43\\
&=(536+243)\div 100\\
&=779\div 100\\
&=7.79\end{aligned}}
\Problem{（2）$9.87-4.52=5.35$}

\MethodTitle{方法一：按题意列竖式}
\Answer{\begin{tikzpicture}[x=1em,y=1em,baseline=(current bounding box.center),draw=answerred,text=answerred]
\node[anchor=base,inner sep=0pt] at (0.000,0) {$9$};
\node[anchor=base,inner sep=0pt] at (0.800,0) {$.$};
\node[anchor=base,inner sep=0pt] at (1.100,0) {$8$};
\node[anchor=base,inner sep=0pt] at (1.900,0) {$7$};
\node[anchor=base,inner sep=0pt] at (0.000,-1.1) {$4$};
\node[anchor=base,inner sep=0pt] at (0.800,-1.1) {$.$};
\node[anchor=base,inner sep=0pt] at (1.100,-1.1) {$5$};
\node[anchor=base,inner sep=0pt] at (1.900,-1.1) {$2$};
\node[anchor=base,inner sep=0pt] at (0.000,-2.3) {$5$};
\node[anchor=base,inner sep=0pt] at (0.800,-2.3) {$.$};
\node[anchor=base,inner sep=0pt] at (1.100,-2.3) {$3$};
\node[anchor=base,inner sep=0pt] at (1.900,-2.3) {$5$};
\node[anchor=base,inner sep=0pt] at (-.7,-1.1) {$-$};
\draw[line width=.5pt] (-1,-1.52)--(2.500,-1.52);
\end{tikzpicture}}

\MethodBlock{方法二：拆成整数部分和小数部分}{上方竖式及题目要求的验算仍保留；本法用于核对。}{\begin{aligned}[t]&\phantom{={}}9.87-4.52\\
&=9.87-4-0.52\\
&=5.87-0.52\\
&=5.35\end{aligned}}

\MethodBlock{方法三：按相同计数单位计算}{上方竖式及题目要求的验算仍保留；本法用于核对。先看成若干个百分之一，按整数加减，最后还原成小数。}{\begin{aligned}[t]&\phantom{={}}9.87-4.52\\
&=(987-452)\div 100\\
&=535\div 100\\
&=5.35\end{aligned}}
\Problem{（3）$6.7+3.28=9.98$}

\MethodTitle{方法一：按题意列竖式}
\Answer{\begin{tikzpicture}[x=1em,y=1em,baseline=(current bounding box.center),draw=answerred,text=answerred]
\node[anchor=base,inner sep=0pt] at (0.000,0) {$6$};
\node[anchor=base,inner sep=0pt] at (0.800,0) {$.$};
\node[anchor=base,inner sep=0pt] at (1.100,0) {$7$};
\node[anchor=base,inner sep=0pt] at (1.900,0) {$0$};
\node[anchor=base,inner sep=0pt] at (0.000,-1.1) {$3$};
\node[anchor=base,inner sep=0pt] at (0.800,-1.1) {$.$};
\node[anchor=base,inner sep=0pt] at (1.100,-1.1) {$2$};
\node[anchor=base,inner sep=0pt] at (1.900,-1.1) {$8$};
\node[anchor=base,inner sep=0pt] at (0.000,-2.3) {$9$};
\node[anchor=base,inner sep=0pt] at (0.800,-2.3) {$.$};
\node[anchor=base,inner sep=0pt] at (1.100,-2.3) {$9$};
\node[anchor=base,inner sep=0pt] at (1.900,-2.3) {$8$};
\node[anchor=base,inner sep=0pt] at (-.7,-1.1) {$+$};
\draw[line width=.5pt] (-1,-1.52)--(2.500,-1.52);
\end{tikzpicture}}

\MethodBlock{方法二：拆成整数部分和小数部分}{上方竖式及题目要求的验算仍保留；本法用于核对。}{\begin{aligned}[t]&\phantom{={}}6.7+3.28\\
&=6.7+3+0.28\\
&=9.7+0.28\\
&=9.98\end{aligned}}

\MethodBlock{方法三：按相同计数单位计算}{上方竖式及题目要求的验算仍保留；本法用于核对。先看成若干个百分之一，按整数加减，最后还原成小数。}{\begin{aligned}[t]&\phantom{={}}6.7+3.28\\
&=(670+328)\div 100\\
&=998\div 100\\
&=9.98\end{aligned}}
\Problem{（4）$15.6-7.09=8.51$}

\MethodTitle{方法一：按题意列竖式}
\Answer{\begin{tikzpicture}[x=1em,y=1em,baseline=(current bounding box.center),draw=answerred,text=answerred]
\node[anchor=base,inner sep=0pt] at (0.000,0) {$1$};
\node[anchor=base,inner sep=0pt] at (0.800,0) {$5$};
\node[anchor=base,inner sep=0pt] at (1.600,0) {$.$};
\node[anchor=base,inner sep=0pt] at (1.900,0) {$6$};
\node[anchor=base,inner sep=0pt] at (2.700,0) {$0$};
\node[anchor=base,inner sep=0pt] at (0.800,-1.1) {$7$};
\node[anchor=base,inner sep=0pt] at (1.600,-1.1) {$.$};
\node[anchor=base,inner sep=0pt] at (1.900,-1.1) {$0$};
\node[anchor=base,inner sep=0pt] at (2.700,-1.1) {$9$};
\node[anchor=base,inner sep=0pt] at (0.800,-2.3) {$8$};
\node[anchor=base,inner sep=0pt] at (1.600,-2.3) {$.$};
\node[anchor=base,inner sep=0pt] at (1.900,-2.3) {$5$};
\node[anchor=base,inner sep=0pt] at (2.700,-2.3) {$1$};
\node[anchor=base,inner sep=0pt] at (-.7,-1.1) {$-$};
\draw[line width=.5pt] (-1,-1.52)--(3.300,-1.52);
\draw[line width=.35pt] (-0.160,.03)--(0.160,.43);
\node[anchor=base,inner sep=0pt,font=\normalsize] at (0.000,.72) {$0$};
\draw[line width=.35pt] (0.640,.03)--(0.960,.43);
\node[anchor=base,inner sep=0pt,font=\normalsize] at (0.800,.72) {$15$};
\draw[line width=.35pt] (1.740,.03)--(2.060,.43);
\node[anchor=base,inner sep=0pt,font=\normalsize] at (1.900,.72) {$5$};
\draw[line width=.35pt] (2.540,.03)--(2.860,.43);
\node[anchor=base,inner sep=0pt,font=\normalsize] at (2.700,.72) {$10$};
\end{tikzpicture}}

\MethodBlock{方法二：拆成整数部分和小数部分}{上方竖式及题目要求的验算仍保留；本法用于核对。}{\begin{aligned}[t]&\phantom{={}}15.6-7.09\\
&=15.6-7-0.09\\
&=8.6-0.09\\
&=8.51\end{aligned}}

\MethodBlock{方法三：按相同计数单位计算}{上方竖式及题目要求的验算仍保留；本法用于核对。先看成若干个百分之一，按整数加减，最后还原成小数。}{\begin{aligned}[t]&\phantom{={}}15.6-7.09\\
&=(1560-709)\div 100\\
&=851\div 100\\
&=8.51\end{aligned}}
\Problem{（5）$7.89+5.64=13.53$}

\MethodTitle{方法一：按题意列竖式}
\Answer{\begin{tikzpicture}[x=1em,y=1em,baseline=(current bounding box.center),draw=answerred,text=answerred]
\node[anchor=base,inner sep=0pt] at (0.800,0) {$7$};
\node[anchor=base,inner sep=0pt] at (1.600,0) {$.$};
\node[anchor=base,inner sep=0pt] at (1.900,0) {$8$};
\node[anchor=base,inner sep=0pt] at (2.700,0) {$9$};
\node[anchor=base,inner sep=0pt] at (0.800,-1.1) {$5$};
\node[anchor=base,inner sep=0pt] at (1.600,-1.1) {$.$};
\node[anchor=base,inner sep=0pt] at (1.900,-1.1) {$6$};
\node[anchor=base,inner sep=0pt] at (2.700,-1.1) {$4$};
\node[anchor=base,inner sep=0pt] at (0.000,-2.3) {$1$};
\node[anchor=base,inner sep=0pt] at (0.800,-2.3) {$3$};
\node[anchor=base,inner sep=0pt] at (1.600,-2.3) {$.$};
\node[anchor=base,inner sep=0pt] at (1.900,-2.3) {$5$};
\node[anchor=base,inner sep=0pt] at (2.700,-2.3) {$3$};
\node[anchor=base,inner sep=0pt] at (-.7,-1.1) {$+$};
\draw[line width=.5pt] (-1,-1.52)--(3.300,-1.52);
\node[anchor=base,inner sep=0pt,font=\normalsize] at (1.900,.72) {$1$};
\node[anchor=base,inner sep=0pt,font=\normalsize] at (0.800,.72) {$1$};
\node[anchor=base,inner sep=0pt,font=\normalsize] at (0.000,.72) {$1$};
\end{tikzpicture}}

\MethodBlock{方法二：拆成整数部分和小数部分}{上方竖式及题目要求的验算仍保留；本法用于核对。}{\begin{aligned}[t]&\phantom{={}}7.89+5.64\\
&=7.89+5+0.64\\
&=12.89+0.64\\
&=13.53\end{aligned}}

\MethodBlock{方法三：按相同计数单位计算}{上方竖式及题目要求的验算仍保留；本法用于核对。先看成若干个百分之一，按整数加减，最后还原成小数。}{\begin{aligned}[t]&\phantom{={}}7.89+5.64\\
&=(789+564)\div 100\\
&=1353\div 100\\
&=13.53\end{aligned}}
\Problem{（6）$12.3-4.75=7.55$}

\MethodTitle{方法一：按题意列竖式}
\Answer{\begin{tikzpicture}[x=1em,y=1em,baseline=(current bounding box.center),draw=answerred,text=answerred]
\node[anchor=base,inner sep=0pt] at (0.000,0) {$1$};
\node[anchor=base,inner sep=0pt] at (0.800,0) {$2$};
\node[anchor=base,inner sep=0pt] at (1.600,0) {$.$};
\node[anchor=base,inner sep=0pt] at (1.900,0) {$3$};
\node[anchor=base,inner sep=0pt] at (2.700,0) {$0$};
\node[anchor=base,inner sep=0pt] at (0.800,-1.1) {$4$};
\node[anchor=base,inner sep=0pt] at (1.600,-1.1) {$.$};
\node[anchor=base,inner sep=0pt] at (1.900,-1.1) {$7$};
\node[anchor=base,inner sep=0pt] at (2.700,-1.1) {$5$};
\node[anchor=base,inner sep=0pt] at (0.800,-2.3) {$7$};
\node[anchor=base,inner sep=0pt] at (1.600,-2.3) {$.$};
\node[anchor=base,inner sep=0pt] at (1.900,-2.3) {$5$};
\node[anchor=base,inner sep=0pt] at (2.700,-2.3) {$5$};
\node[anchor=base,inner sep=0pt] at (-.7,-1.1) {$-$};
\draw[line width=.5pt] (-1,-1.52)--(3.300,-1.52);
\draw[line width=.35pt] (-0.160,.03)--(0.160,.43);
\node[anchor=base,inner sep=0pt,font=\normalsize] at (0.000,.72) {$0$};
\draw[line width=.35pt] (0.640,.03)--(0.960,.43);
\node[anchor=base,inner sep=0pt,font=\normalsize] at (0.800,.72) {$11$};
\draw[line width=.35pt] (1.740,.03)--(2.060,.43);
\node[anchor=base,inner sep=0pt,font=\normalsize] at (1.900,.72) {$12$};
\draw[line width=.35pt] (2.540,.03)--(2.860,.43);
\node[anchor=base,inner sep=0pt,font=\normalsize] at (2.700,.72) {$10$};
\end{tikzpicture}}

\MethodBlock{方法二：多减再加回}{上方竖式及题目要求的验算仍保留；本法用于核对。减去凑成的整数后，把多减的部分加回来。}{\begin{aligned}[t]&\phantom{={}}12.3-4.75\\
&=12.3-5+0.25\\
&=7.3+0.25\\
&=7.55\end{aligned}}

\MethodBlock{方法三：拆成整数部分和小数部分}{上方竖式及题目要求的验算仍保留；本法用于核对。}{\begin{aligned}[t]&\phantom{={}}12.3-4.75\\
&=12.3-4-0.75\\
&=8.3-0.75\\
&=7.55\end{aligned}}
\Problem{（7）$0.85+4.67=5.52$}

\MethodTitle{方法一：按题意列竖式}
\Answer{\begin{tikzpicture}[x=1em,y=1em,baseline=(current bounding box.center),draw=answerred,text=answerred]
\node[anchor=base,inner sep=0pt] at (0.000,0) {$0$};
\node[anchor=base,inner sep=0pt] at (0.800,0) {$.$};
\node[anchor=base,inner sep=0pt] at (1.100,0) {$8$};
\node[anchor=base,inner sep=0pt] at (1.900,0) {$5$};
\node[anchor=base,inner sep=0pt] at (0.000,-1.1) {$4$};
\node[anchor=base,inner sep=0pt] at (0.800,-1.1) {$.$};
\node[anchor=base,inner sep=0pt] at (1.100,-1.1) {$6$};
\node[anchor=base,inner sep=0pt] at (1.900,-1.1) {$7$};
\node[anchor=base,inner sep=0pt] at (0.000,-2.3) {$5$};
\node[anchor=base,inner sep=0pt] at (0.800,-2.3) {$.$};
\node[anchor=base,inner sep=0pt] at (1.100,-2.3) {$5$};
\node[anchor=base,inner sep=0pt] at (1.900,-2.3) {$2$};
\node[anchor=base,inner sep=0pt] at (-.7,-1.1) {$+$};
\draw[line width=.5pt] (-1,-1.52)--(2.500,-1.52);
\node[anchor=base,inner sep=0pt,font=\normalsize] at (1.100,.72) {$1$};
\node[anchor=base,inner sep=0pt,font=\normalsize] at (0.000,.72) {$1$};
\end{tikzpicture}}

\MethodBlock{方法二：拆成整数部分和小数部分}{上方竖式及题目要求的验算仍保留；本法用于核对。}{\begin{aligned}[t]&\phantom{={}}0.85+4.67\\
&=0.85+4+0.67\\
&=4.85+0.67\\
&=5.52\end{aligned}}

\MethodBlock{方法三：按相同计数单位计算}{上方竖式及题目要求的验算仍保留；本法用于核对。先看成若干个百分之一，按整数加减，最后还原成小数。}{\begin{aligned}[t]&\phantom{={}}0.85+4.67\\
&=(85+467)\div 100\\
&=552\div 100\\
&=5.52\end{aligned}}
\Problem{（8）$20-9.36=10.64$}

\MethodTitle{方法一：按题意列竖式}
\Answer{\begin{tikzpicture}[x=1em,y=1em,baseline=(current bounding box.center),draw=answerred,text=answerred]
\node[anchor=base,inner sep=0pt] at (0.000,0) {$2$};
\node[anchor=base,inner sep=0pt] at (0.800,0) {$0$};
\node[anchor=base,inner sep=0pt] at (1.600,0) {$.$};
\node[anchor=base,inner sep=0pt] at (1.900,0) {$0$};
\node[anchor=base,inner sep=0pt] at (2.700,0) {$0$};
\node[anchor=base,inner sep=0pt] at (0.800,-1.1) {$9$};
\node[anchor=base,inner sep=0pt] at (1.600,-1.1) {$.$};
\node[anchor=base,inner sep=0pt] at (1.900,-1.1) {$3$};
\node[anchor=base,inner sep=0pt] at (2.700,-1.1) {$6$};
\node[anchor=base,inner sep=0pt] at (0.000,-2.3) {$1$};
\node[anchor=base,inner sep=0pt] at (0.800,-2.3) {$0$};
\node[anchor=base,inner sep=0pt] at (1.600,-2.3) {$.$};
\node[anchor=base,inner sep=0pt] at (1.900,-2.3) {$6$};
\node[anchor=base,inner sep=0pt] at (2.700,-2.3) {$4$};
\node[anchor=base,inner sep=0pt] at (-.7,-1.1) {$-$};
\draw[line width=.5pt] (-1,-1.52)--(3.300,-1.52);
\draw[line width=.35pt] (-0.160,.03)--(0.160,.43);
\node[anchor=base,inner sep=0pt,font=\normalsize] at (0.000,.72) {$1$};
\draw[line width=.35pt] (0.640,.03)--(0.960,.43);
\node[anchor=base,inner sep=0pt,font=\normalsize] at (0.800,.72) {$9$};
\draw[line width=.35pt] (1.740,.03)--(2.060,.43);
\node[anchor=base,inner sep=0pt,font=\normalsize] at (1.900,.72) {$9$};
\draw[line width=.35pt] (2.540,.03)--(2.860,.43);
\node[anchor=base,inner sep=0pt,font=\normalsize] at (2.700,.72) {$10$};
\end{tikzpicture}}

\MethodBlock{方法二：拆成整数部分和小数部分}{上方竖式及题目要求的验算仍保留；本法用于核对。}{\begin{aligned}[t]&\phantom{={}}20-9.36\\
&=20-9-0.36\\
&=11-0.36\\
&=10.64\end{aligned}}

\MethodBlock{方法三：按相同计数单位计算}{上方竖式及题目要求的验算仍保留；本法用于核对。先看成若干个百分之一，按整数加减，最后还原成小数。}{\begin{aligned}[t]&\phantom{={}}20-9.36\\
&=(2000-936)\div 100\\
&=1064\div 100\\
&=10.64\end{aligned}}
\Problem{（9）$8.35+4.65=13$}

\MethodTitle{方法一：按题意列竖式}
\Answer{\begin{tikzpicture}[x=1em,y=1em,baseline=(current bounding box.center),draw=answerred,text=answerred]
\node[anchor=base,inner sep=0pt] at (0.800,0) {$8$};
\node[anchor=base,inner sep=0pt] at (1.600,0) {$.$};
\node[anchor=base,inner sep=0pt] at (1.900,0) {$3$};
\node[anchor=base,inner sep=0pt] at (2.700,0) {$5$};
\node[anchor=base,inner sep=0pt] at (0.800,-1.1) {$4$};
\node[anchor=base,inner sep=0pt] at (1.600,-1.1) {$.$};
\node[anchor=base,inner sep=0pt] at (1.900,-1.1) {$6$};
\node[anchor=base,inner sep=0pt] at (2.700,-1.1) {$5$};
\node[anchor=base,inner sep=0pt] at (0.000,-2.3) {$1$};
\node[anchor=base,inner sep=0pt] at (0.800,-2.3) {$3$};
\node[anchor=base,inner sep=0pt] at (1.600,-2.3) {$.$};
\node[anchor=base,inner sep=0pt] at (1.900,-2.3) {$0$};
\node[anchor=base,inner sep=0pt] at (2.700,-2.3) {$0$};
\node[anchor=base,inner sep=0pt] at (-.7,-1.1) {$+$};
\draw[line width=.5pt] (-1,-1.52)--(3.300,-1.52);
\node[anchor=base,inner sep=0pt,font=\normalsize] at (1.900,.72) {$1$};
\node[anchor=base,inner sep=0pt,font=\normalsize] at (0.800,.72) {$1$};
\node[anchor=base,inner sep=0pt,font=\normalsize] at (0.000,.72) {$1$};
\end{tikzpicture}}

\MethodBlock{方法二：拆成整数部分和小数部分}{上方竖式及题目要求的验算仍保留；本法用于核对。}{\begin{aligned}[t]&\phantom{={}}8.35+4.65\\
&=8.35+4+0.65\\
&=12.35+0.65\\
&=13\end{aligned}}

\MethodBlock{方法三：按相同计数单位计算}{上方竖式及题目要求的验算仍保留；本法用于核对。先看成若干个百分之一，按整数加减，最后还原成小数。}{\begin{aligned}[t]&\phantom{={}}8.35+4.65\\
&=(835+465)\div 100\\
&=1300\div 100\\
&=13\end{aligned}}
\Problem{（10）$1.19+7.9=9.09$}

\MethodTitle{方法一：按题意列竖式}
\Answer{\begin{tikzpicture}[x=1em,y=1em,baseline=(current bounding box.center),draw=answerred,text=answerred]
\node[anchor=base,inner sep=0pt] at (0.000,0) {$1$};
\node[anchor=base,inner sep=0pt] at (0.800,0) {$.$};
\node[anchor=base,inner sep=0pt] at (1.100,0) {$1$};
\node[anchor=base,inner sep=0pt] at (1.900,0) {$9$};
\node[anchor=base,inner sep=0pt] at (0.000,-1.1) {$7$};
\node[anchor=base,inner sep=0pt] at (0.800,-1.1) {$.$};
\node[anchor=base,inner sep=0pt] at (1.100,-1.1) {$9$};
\node[anchor=base,inner sep=0pt] at (1.900,-1.1) {$0$};
\node[anchor=base,inner sep=0pt] at (0.000,-2.3) {$9$};
\node[anchor=base,inner sep=0pt] at (0.800,-2.3) {$.$};
\node[anchor=base,inner sep=0pt] at (1.100,-2.3) {$0$};
\node[anchor=base,inner sep=0pt] at (1.900,-2.3) {$9$};
\node[anchor=base,inner sep=0pt] at (-.7,-1.1) {$+$};
\draw[line width=.5pt] (-1,-1.52)--(2.500,-1.52);
\node[anchor=base,inner sep=0pt,font=\normalsize] at (0.000,.72) {$1$};
\end{tikzpicture}}

\MethodBlock{方法二：凑整再减去多加的数}{上方竖式及题目要求的验算仍保留；本法用于核对。把第二个数凑成整数，再减去多加的部分。}{\begin{aligned}[t]&\phantom{={}}1.19+7.9\\
&=1.19+8-0.1\\
&=9.19-0.1\\
&=9.09\end{aligned}}

\MethodBlock{方法三：拆成整数部分和小数部分}{上方竖式及题目要求的验算仍保留；本法用于核对。}{\begin{aligned}[t]&\phantom{={}}1.19+7.9\\
&=1.19+7+0.9\\
&=8.19+0.9\\
&=9.09\end{aligned}}
\Problem{（11）$10-3.28=6.72$}

\MethodTitle{方法一：按题意列竖式}
\Answer{\begin{tikzpicture}[x=1em,y=1em,baseline=(current bounding box.center),draw=answerred,text=answerred]
\node[anchor=base,inner sep=0pt] at (0.000,0) {$1$};
\node[anchor=base,inner sep=0pt] at (0.800,0) {$0$};
\node[anchor=base,inner sep=0pt] at (1.600,0) {$.$};
\node[anchor=base,inner sep=0pt] at (1.900,0) {$0$};
\node[anchor=base,inner sep=0pt] at (2.700,0) {$0$};
\node[anchor=base,inner sep=0pt] at (0.800,-1.1) {$3$};
\node[anchor=base,inner sep=0pt] at (1.600,-1.1) {$.$};
\node[anchor=base,inner sep=0pt] at (1.900,-1.1) {$2$};
\node[anchor=base,inner sep=0pt] at (2.700,-1.1) {$8$};
\node[anchor=base,inner sep=0pt] at (0.800,-2.3) {$6$};
\node[anchor=base,inner sep=0pt] at (1.600,-2.3) {$.$};
\node[anchor=base,inner sep=0pt] at (1.900,-2.3) {$7$};
\node[anchor=base,inner sep=0pt] at (2.700,-2.3) {$2$};
\node[anchor=base,inner sep=0pt] at (-.7,-1.1) {$-$};
\draw[line width=.5pt] (-1,-1.52)--(3.300,-1.52);
\draw[line width=.35pt] (-0.160,.03)--(0.160,.43);
\node[anchor=base,inner sep=0pt,font=\normalsize] at (0.000,.72) {$0$};
\draw[line width=.35pt] (0.640,.03)--(0.960,.43);
\node[anchor=base,inner sep=0pt,font=\normalsize] at (0.800,.72) {$9$};
\draw[line width=.35pt] (1.740,.03)--(2.060,.43);
\node[anchor=base,inner sep=0pt,font=\normalsize] at (1.900,.72) {$9$};
\draw[line width=.35pt] (2.540,.03)--(2.860,.43);
\node[anchor=base,inner sep=0pt,font=\normalsize] at (2.700,.72) {$10$};
\end{tikzpicture}}

\MethodBlock{方法二：拆成整数部分和小数部分}{上方竖式及题目要求的验算仍保留；本法用于核对。}{\begin{aligned}[t]&\phantom{={}}10-3.28\\
&=10-3-0.28\\
&=7-0.28\\
&=6.72\end{aligned}}

\MethodBlock{方法三：按相同计数单位计算}{上方竖式及题目要求的验算仍保留；本法用于核对。先看成若干个百分之一，按整数加减，最后还原成小数。}{\begin{aligned}[t]&\phantom{={}}10-3.28\\
&=(1000-328)\div 100\\
&=672\div 100\\
&=6.72\end{aligned}}
\Problem{（12）$15.3-9.56=5.74$}

\MethodTitle{方法一：按题意列竖式}
\Answer{\begin{tikzpicture}[x=1em,y=1em,baseline=(current bounding box.center),draw=answerred,text=answerred]
\node[anchor=base,inner sep=0pt] at (0.000,0) {$1$};
\node[anchor=base,inner sep=0pt] at (0.800,0) {$5$};
\node[anchor=base,inner sep=0pt] at (1.600,0) {$.$};
\node[anchor=base,inner sep=0pt] at (1.900,0) {$3$};
\node[anchor=base,inner sep=0pt] at (2.700,0) {$0$};
\node[anchor=base,inner sep=0pt] at (0.800,-1.1) {$9$};
\node[anchor=base,inner sep=0pt] at (1.600,-1.1) {$.$};
\node[anchor=base,inner sep=0pt] at (1.900,-1.1) {$5$};
\node[anchor=base,inner sep=0pt] at (2.700,-1.1) {$6$};
\node[anchor=base,inner sep=0pt] at (0.800,-2.3) {$5$};
\node[anchor=base,inner sep=0pt] at (1.600,-2.3) {$.$};
\node[anchor=base,inner sep=0pt] at (1.900,-2.3) {$7$};
\node[anchor=base,inner sep=0pt] at (2.700,-2.3) {$4$};
\node[anchor=base,inner sep=0pt] at (-.7,-1.1) {$-$};
\draw[line width=.5pt] (-1,-1.52)--(3.300,-1.52);
\draw[line width=.35pt] (-0.160,.03)--(0.160,.43);
\node[anchor=base,inner sep=0pt,font=\normalsize] at (0.000,.72) {$0$};
\draw[line width=.35pt] (0.640,.03)--(0.960,.43);
\node[anchor=base,inner sep=0pt,font=\normalsize] at (0.800,.72) {$14$};
\draw[line width=.35pt] (1.740,.03)--(2.060,.43);
\node[anchor=base,inner sep=0pt,font=\normalsize] at (1.900,.72) {$12$};
\draw[line width=.35pt] (2.540,.03)--(2.860,.43);
\node[anchor=base,inner sep=0pt,font=\normalsize] at (2.700,.72) {$10$};
\end{tikzpicture}}

\MethodBlock{方法二：拆成整数部分和小数部分}{上方竖式及题目要求的验算仍保留；本法用于核对。}{\begin{aligned}[t]&\phantom{={}}15.3-9.56\\
&=15.3-9-0.56\\
&=6.3-0.56\\
&=5.74\end{aligned}}

\MethodBlock{方法三：按相同计数单位计算}{上方竖式及题目要求的验算仍保留；本法用于核对。先看成若干个百分之一，按整数加减，最后还原成小数。}{\begin{aligned}[t]&\phantom{={}}15.3-9.56\\
&=(1530-956)\div 100\\
&=574\div 100\\
&=5.74\end{aligned}}
\par\noindent

\Answer{
\QHeading{计算过程：2. 列竖式计算}

小数点对齐，缺位补零；从右往左逐列计算。下面完整列出进位、退位及各数位的计算过程。

\par\medskip\noindent
\textbf{第1小题：\(5.36+2.43=7.79\)}\par\smallskip
\noindent
\begin{tikzpicture}[x=1em,y=1em,baseline=(current bounding box.center),draw=answerred,text=answerred]
\node[anchor=base,inner sep=0pt] at (0.000,0) {$5$};
\node[anchor=base,inner sep=0pt] at (0.800,0) {$.$};
\node[anchor=base,inner sep=0pt] at (1.100,0) {$3$};
\node[anchor=base,inner sep=0pt] at (1.900,0) {$6$};
\node[anchor=base,inner sep=0pt] at (0.000,-1.1) {$2$};
\node[anchor=base,inner sep=0pt] at (0.800,-1.1) {$.$};
\node[anchor=base,inner sep=0pt] at (1.100,-1.1) {$4$};
\node[anchor=base,inner sep=0pt] at (1.900,-1.1) {$3$};
\node[anchor=base,inner sep=0pt] at (0.000,-2.3) {$7$};
\node[anchor=base,inner sep=0pt] at (0.800,-2.3) {$.$};
\node[anchor=base,inner sep=0pt] at (1.100,-2.3) {$7$};
\node[anchor=base,inner sep=0pt] at (1.900,-2.3) {$9$};
\node[anchor=base,inner sep=0pt] at (-.7,-1.1) {$+$};
\draw[line width=.5pt] (-1,-1.52)--(2.500,-1.52);
\end{tikzpicture}
\par\smallskip{\normalsize 上方小字表示从后一位进来的数。\par}
\vspace{0pt}
\begin{enumerate}[leftmargin=1.7em,itemsep=5pt,topsep=0pt]
\item 百分位：\(\begin{aligned}[t]&\phantom{={}}6+3\\&=9\end{aligned}\)\quad 本位写9。
\item 十分位：\(\begin{aligned}[t]&\phantom{={}}3+4\\&=7\end{aligned}\)\quad 本位写7。
\item 个位：\(\begin{aligned}[t]&\phantom{={}}5+2\\&=7\end{aligned}\)\quad 本位写7。
\end{enumerate}
\par\smallskip
因此，\(\begin{aligned}[t]&\phantom{={}}5.36+2.43\\&=7.79\end{aligned}\)。
\par\medskip

\par\medskip\noindent
\textbf{第2小题：\(9.87-4.52=5.35\)}\par\smallskip
\noindent
\begin{tikzpicture}[x=1em,y=1em,baseline=(current bounding box.center),draw=answerred,text=answerred]
\node[anchor=base,inner sep=0pt] at (0.000,0) {$9$};
\node[anchor=base,inner sep=0pt] at (0.800,0) {$.$};
\node[anchor=base,inner sep=0pt] at (1.100,0) {$8$};
\node[anchor=base,inner sep=0pt] at (1.900,0) {$7$};
\node[anchor=base,inner sep=0pt] at (0.000,-1.1) {$4$};
\node[anchor=base,inner sep=0pt] at (0.800,-1.1) {$.$};
\node[anchor=base,inner sep=0pt] at (1.100,-1.1) {$5$};
\node[anchor=base,inner sep=0pt] at (1.900,-1.1) {$2$};
\node[anchor=base,inner sep=0pt] at (0.000,-2.3) {$5$};
\node[anchor=base,inner sep=0pt] at (0.800,-2.3) {$.$};
\node[anchor=base,inner sep=0pt] at (1.100,-2.3) {$3$};
\node[anchor=base,inner sep=0pt] at (1.900,-2.3) {$5$};
\node[anchor=base,inner sep=0pt] at (-.7,-1.1) {$-$};
\draw[line width=.5pt] (-1,-1.52)--(2.500,-1.52);
\end{tikzpicture}
\par\smallskip{\normalsize 划线数字是退位前的数，上方小字是退位后参与计算的数。\par}
\vspace{0pt}
\begin{enumerate}[leftmargin=1.7em,itemsep=5pt,topsep=0pt]
\item 百分位：\(\begin{aligned}[t]&\phantom{={}}7-2\\&=5\end{aligned}\)\quad 本位写5。
\item 十分位：\(\begin{aligned}[t]&\phantom{={}}8-5\\&=3\end{aligned}\)\quad 本位写3。
\item 个位：\(\begin{aligned}[t]&\phantom{={}}9-4\\&=5\end{aligned}\)\quad 本位写5。
\end{enumerate}
\par\smallskip
因此，\(\begin{aligned}[t]&\phantom{={}}9.87-4.52\\&=5.35\end{aligned}\)。
\par\medskip

\par\medskip\noindent
\textbf{第3小题：\(6.7+3.28=9.98\)}\par\smallskip
\noindent
\begin{tikzpicture}[x=1em,y=1em,baseline=(current bounding box.center),draw=answerred,text=answerred]
\node[anchor=base,inner sep=0pt] at (0.000,0) {$6$};
\node[anchor=base,inner sep=0pt] at (0.800,0) {$.$};
\node[anchor=base,inner sep=0pt] at (1.100,0) {$7$};
\node[anchor=base,inner sep=0pt] at (1.900,0) {$0$};
\node[anchor=base,inner sep=0pt] at (0.000,-1.1) {$3$};
\node[anchor=base,inner sep=0pt] at (0.800,-1.1) {$.$};
\node[anchor=base,inner sep=0pt] at (1.100,-1.1) {$2$};
\node[anchor=base,inner sep=0pt] at (1.900,-1.1) {$8$};
\node[anchor=base,inner sep=0pt] at (0.000,-2.3) {$9$};
\node[anchor=base,inner sep=0pt] at (0.800,-2.3) {$.$};
\node[anchor=base,inner sep=0pt] at (1.100,-2.3) {$9$};
\node[anchor=base,inner sep=0pt] at (1.900,-2.3) {$8$};
\node[anchor=base,inner sep=0pt] at (-.7,-1.1) {$+$};
\draw[line width=.5pt] (-1,-1.52)--(2.500,-1.52);
\end{tikzpicture}
\par\smallskip{\normalsize 上方小字表示从后一位进来的数。\par}
\vspace{0pt}
\begin{enumerate}[leftmargin=1.7em,itemsep=5pt,topsep=0pt]
\item 百分位：\(\begin{aligned}[t]&\phantom{={}}0+8\\&=8\end{aligned}\)\quad 本位写8。
\item 十分位：\(\begin{aligned}[t]&\phantom{={}}7+2\\&=9\end{aligned}\)\quad 本位写9。
\item 个位：\(\begin{aligned}[t]&\phantom{={}}6+3\\&=9\end{aligned}\)\quad 本位写9。
\end{enumerate}
\par\smallskip
因此，\(\begin{aligned}[t]&\phantom{={}}6.7+3.28\\&=9.98\end{aligned}\)。
\par\medskip

\par\medskip\noindent
\textbf{第4小题：\(15.6-7.09=8.51\)}\par\smallskip
\noindent
\begin{tikzpicture}[x=1em,y=1em,baseline=(current bounding box.center),draw=answerred,text=answerred]
\node[anchor=base,inner sep=0pt] at (0.000,0) {$1$};
\node[anchor=base,inner sep=0pt] at (0.800,0) {$5$};
\node[anchor=base,inner sep=0pt] at (1.600,0) {$.$};
\node[anchor=base,inner sep=0pt] at (1.900,0) {$6$};
\node[anchor=base,inner sep=0pt] at (2.700,0) {$0$};
\node[anchor=base,inner sep=0pt] at (0.800,-1.1) {$7$};
\node[anchor=base,inner sep=0pt] at (1.600,-1.1) {$.$};
\node[anchor=base,inner sep=0pt] at (1.900,-1.1) {$0$};
\node[anchor=base,inner sep=0pt] at (2.700,-1.1) {$9$};
\node[anchor=base,inner sep=0pt] at (0.800,-2.3) {$8$};
\node[anchor=base,inner sep=0pt] at (1.600,-2.3) {$.$};
\node[anchor=base,inner sep=0pt] at (1.900,-2.3) {$5$};
\node[anchor=base,inner sep=0pt] at (2.700,-2.3) {$1$};
\node[anchor=base,inner sep=0pt] at (-.7,-1.1) {$-$};
\draw[line width=.5pt] (-1,-1.52)--(3.300,-1.52);
\draw[line width=.35pt] (-0.160,.03)--(0.160,.43);
\node[anchor=base,inner sep=0pt,font=\normalsize] at (0.000,.72) {$0$};
\draw[line width=.35pt] (0.640,.03)--(0.960,.43);
\node[anchor=base,inner sep=0pt,font=\normalsize] at (0.800,.72) {$15$};
\draw[line width=.35pt] (1.740,.03)--(2.060,.43);
\node[anchor=base,inner sep=0pt,font=\normalsize] at (1.900,.72) {$5$};
\draw[line width=.35pt] (2.540,.03)--(2.860,.43);
\node[anchor=base,inner sep=0pt,font=\normalsize] at (2.700,.72) {$10$};
\end{tikzpicture}
\par\smallskip{\normalsize 划线数字是退位前的数，上方小字是退位后参与计算的数。\par}
\vspace{0pt}
\begin{enumerate}[leftmargin=1.7em,itemsep=5pt,topsep=0pt]
\item 十分位退位：\(\begin{aligned}[t]&\phantom{={}}6-1\\&=5\end{aligned}\)\quad 借出1给后一位。
\item 百分位接收退位：\(\begin{aligned}[t]&\phantom{={}}0+10\\&=10\end{aligned}\)\quad 前一位的1变成本位的10。
\item 百分位：\(\begin{aligned}[t]&\phantom{={}}10-9\\&=1\end{aligned}\)\quad 本位写1。
\item 十分位：\(\begin{aligned}[t]&\phantom{={}}5-0\\&=5\end{aligned}\)\quad 本位写5。
\item 十位退位：\(\begin{aligned}[t]&\phantom{={}}1-1\\&=0\end{aligned}\)\quad 借出1给后一位。
\item 个位接收退位：\(\begin{aligned}[t]&\phantom{={}}5+10\\&=15\end{aligned}\)\quad 前一位的1变成本位的10。
\item 个位：\(\begin{aligned}[t]&\phantom{={}}15-7\\&=8\end{aligned}\)\quad 本位写8。
\item 十位：\(\begin{aligned}[t]&\phantom{={}}0-0\\&=0\end{aligned}\)\quad 本位写0。
\end{enumerate}
\par\smallskip
因此，\(\begin{aligned}[t]&\phantom{={}}15.6-7.09\\&=8.51\end{aligned}\)。
\par\medskip

\par\medskip\noindent
\textbf{第5小题：\(7.89+5.64=13.53\)}\par\smallskip
\noindent
\begin{tikzpicture}[x=1em,y=1em,baseline=(current bounding box.center),draw=answerred,text=answerred]
\node[anchor=base,inner sep=0pt] at (0.800,0) {$7$};
\node[anchor=base,inner sep=0pt] at (1.600,0) {$.$};
\node[anchor=base,inner sep=0pt] at (1.900,0) {$8$};
\node[anchor=base,inner sep=0pt] at (2.700,0) {$9$};
\node[anchor=base,inner sep=0pt] at (0.800,-1.1) {$5$};
\node[anchor=base,inner sep=0pt] at (1.600,-1.1) {$.$};
\node[anchor=base,inner sep=0pt] at (1.900,-1.1) {$6$};
\node[anchor=base,inner sep=0pt] at (2.700,-1.1) {$4$};
\node[anchor=base,inner sep=0pt] at (0.000,-2.3) {$1$};
\node[anchor=base,inner sep=0pt] at (0.800,-2.3) {$3$};
\node[anchor=base,inner sep=0pt] at (1.600,-2.3) {$.$};
\node[anchor=base,inner sep=0pt] at (1.900,-2.3) {$5$};
\node[anchor=base,inner sep=0pt] at (2.700,-2.3) {$3$};
\node[anchor=base,inner sep=0pt] at (-.7,-1.1) {$+$};
\draw[line width=.5pt] (-1,-1.52)--(3.300,-1.52);
\node[anchor=base,inner sep=0pt,font=\normalsize] at (1.900,.72) {$1$};
\node[anchor=base,inner sep=0pt,font=\normalsize] at (0.800,.72) {$1$};
\node[anchor=base,inner sep=0pt,font=\normalsize] at (0.000,.72) {$1$};
\end{tikzpicture}
\par\smallskip{\normalsize 上方小字表示从后一位进来的数。\par}
\vspace{0pt}
\begin{enumerate}[leftmargin=1.7em,itemsep=5pt,topsep=0pt]
\item 百分位：\(\begin{aligned}[t]&\phantom{={}}9+4\\&=13\end{aligned}\)\quad 本位写3，向前一位进1。
\item 十分位：\(\begin{aligned}[t]&\phantom{={}}8+6+1\\&=14+1\\&=15\end{aligned}\)\quad 本位写5，向前一位进1。
\item 个位：\(\begin{aligned}[t]&\phantom{={}}7+5+1\\&=12+1\\&=13\end{aligned}\)\quad 本位写3，向前一位进1。
\end{enumerate}
最高位再写进上来的1。\par
\par\smallskip
因此，\(\begin{aligned}[t]&\phantom{={}}7.89+5.64\\&=13.53\end{aligned}\)。
\par\medskip

\par\medskip\noindent
\textbf{第6小题：\(12.3-4.75=7.55\)}\par\smallskip
\noindent
\begin{tikzpicture}[x=1em,y=1em,baseline=(current bounding box.center),draw=answerred,text=answerred]
\node[anchor=base,inner sep=0pt] at (0.000,0) {$1$};
\node[anchor=base,inner sep=0pt] at (0.800,0) {$2$};
\node[anchor=base,inner sep=0pt] at (1.600,0) {$.$};
\node[anchor=base,inner sep=0pt] at (1.900,0) {$3$};
\node[anchor=base,inner sep=0pt] at (2.700,0) {$0$};
\node[anchor=base,inner sep=0pt] at (0.800,-1.1) {$4$};
\node[anchor=base,inner sep=0pt] at (1.600,-1.1) {$.$};
\node[anchor=base,inner sep=0pt] at (1.900,-1.1) {$7$};
\node[anchor=base,inner sep=0pt] at (2.700,-1.1) {$5$};
\node[anchor=base,inner sep=0pt] at (0.800,-2.3) {$7$};
\node[anchor=base,inner sep=0pt] at (1.600,-2.3) {$.$};
\node[anchor=base,inner sep=0pt] at (1.900,-2.3) {$5$};
\node[anchor=base,inner sep=0pt] at (2.700,-2.3) {$5$};
\node[anchor=base,inner sep=0pt] at (-.7,-1.1) {$-$};
\draw[line width=.5pt] (-1,-1.52)--(3.300,-1.52);
\draw[line width=.35pt] (-0.160,.03)--(0.160,.43);
\node[anchor=base,inner sep=0pt,font=\normalsize] at (0.000,.72) {$0$};
\draw[line width=.35pt] (0.640,.03)--(0.960,.43);
\node[anchor=base,inner sep=0pt,font=\normalsize] at (0.800,.72) {$11$};
\draw[line width=.35pt] (1.740,.03)--(2.060,.43);
\node[anchor=base,inner sep=0pt,font=\normalsize] at (1.900,.72) {$12$};
\draw[line width=.35pt] (2.540,.03)--(2.860,.43);
\node[anchor=base,inner sep=0pt,font=\normalsize] at (2.700,.72) {$10$};
\end{tikzpicture}
\par\smallskip{\normalsize 划线数字是退位前的数，上方小字是退位后参与计算的数。\par}
\vspace{0pt}
\begin{enumerate}[leftmargin=1.7em,itemsep=5pt,topsep=0pt]
\item 十分位退位：\(\begin{aligned}[t]&\phantom{={}}3-1\\&=2\end{aligned}\)\quad 借出1给后一位。
\item 百分位接收退位：\(\begin{aligned}[t]&\phantom{={}}0+10\\&=10\end{aligned}\)\quad 前一位的1变成本位的10。
\item 百分位：\(\begin{aligned}[t]&\phantom{={}}10-5\\&=5\end{aligned}\)\quad 本位写5。
\item 个位退位：\(\begin{aligned}[t]&\phantom{={}}2-1\\&=1\end{aligned}\)\quad 借出1给后一位。
\item 十分位接收退位：\(\begin{aligned}[t]&\phantom{={}}2+10\\&=12\end{aligned}\)\quad 前一位的1变成本位的10。
\item 十分位：\(\begin{aligned}[t]&\phantom{={}}12-7\\&=5\end{aligned}\)\quad 本位写5。
\item 十位退位：\(\begin{aligned}[t]&\phantom{={}}1-1\\&=0\end{aligned}\)\quad 借出1给后一位。
\item 个位接收退位：\(\begin{aligned}[t]&\phantom{={}}1+10\\&=11\end{aligned}\)\quad 前一位的1变成本位的10。
\item 个位：\(\begin{aligned}[t]&\phantom{={}}11-4\\&=7\end{aligned}\)\quad 本位写7。
\item 十位：\(\begin{aligned}[t]&\phantom{={}}0-0\\&=0\end{aligned}\)\quad 本位写0。
\end{enumerate}
\par\smallskip
因此，\(\begin{aligned}[t]&\phantom{={}}12.3-4.75\\&=7.55\end{aligned}\)。
\par\medskip

\par\medskip\noindent
\textbf{第7小题：\(0.85+4.67=5.52\)}\par\smallskip
\noindent
\begin{tikzpicture}[x=1em,y=1em,baseline=(current bounding box.center),draw=answerred,text=answerred]
\node[anchor=base,inner sep=0pt] at (0.000,0) {$0$};
\node[anchor=base,inner sep=0pt] at (0.800,0) {$.$};
\node[anchor=base,inner sep=0pt] at (1.100,0) {$8$};
\node[anchor=base,inner sep=0pt] at (1.900,0) {$5$};
\node[anchor=base,inner sep=0pt] at (0.000,-1.1) {$4$};
\node[anchor=base,inner sep=0pt] at (0.800,-1.1) {$.$};
\node[anchor=base,inner sep=0pt] at (1.100,-1.1) {$6$};
\node[anchor=base,inner sep=0pt] at (1.900,-1.1) {$7$};
\node[anchor=base,inner sep=0pt] at (0.000,-2.3) {$5$};
\node[anchor=base,inner sep=0pt] at (0.800,-2.3) {$.$};
\node[anchor=base,inner sep=0pt] at (1.100,-2.3) {$5$};
\node[anchor=base,inner sep=0pt] at (1.900,-2.3) {$2$};
\node[anchor=base,inner sep=0pt] at (-.7,-1.1) {$+$};
\draw[line width=.5pt] (-1,-1.52)--(2.500,-1.52);
\node[anchor=base,inner sep=0pt,font=\normalsize] at (1.100,.72) {$1$};
\node[anchor=base,inner sep=0pt,font=\normalsize] at (0.000,.72) {$1$};
\end{tikzpicture}
\par\smallskip{\normalsize 上方小字表示从后一位进来的数。\par}
\vspace{0pt}
\begin{enumerate}[leftmargin=1.7em,itemsep=5pt,topsep=0pt]
\item 百分位：\(\begin{aligned}[t]&\phantom{={}}5+7\\&=12\end{aligned}\)\quad 本位写2，向前一位进1。
\item 十分位：\(\begin{aligned}[t]&\phantom{={}}8+6+1\\&=14+1\\&=15\end{aligned}\)\quad 本位写5，向前一位进1。
\item 个位：\(\begin{aligned}[t]&\phantom{={}}0+4+1\\&=4+1\\&=5\end{aligned}\)\quad 本位写5。
\end{enumerate}
\par\smallskip
因此，\(\begin{aligned}[t]&\phantom{={}}0.85+4.67\\&=5.52\end{aligned}\)。
\par\medskip

\par\medskip\noindent
\textbf{第8小题：\(20-9.36=10.64\)}\par\smallskip
\noindent
\begin{tikzpicture}[x=1em,y=1em,baseline=(current bounding box.center),draw=answerred,text=answerred]
\node[anchor=base,inner sep=0pt] at (0.000,0) {$2$};
\node[anchor=base,inner sep=0pt] at (0.800,0) {$0$};
\node[anchor=base,inner sep=0pt] at (1.600,0) {$.$};
\node[anchor=base,inner sep=0pt] at (1.900,0) {$0$};
\node[anchor=base,inner sep=0pt] at (2.700,0) {$0$};
\node[anchor=base,inner sep=0pt] at (0.800,-1.1) {$9$};
\node[anchor=base,inner sep=0pt] at (1.600,-1.1) {$.$};
\node[anchor=base,inner sep=0pt] at (1.900,-1.1) {$3$};
\node[anchor=base,inner sep=0pt] at (2.700,-1.1) {$6$};
\node[anchor=base,inner sep=0pt] at (0.000,-2.3) {$1$};
\node[anchor=base,inner sep=0pt] at (0.800,-2.3) {$0$};
\node[anchor=base,inner sep=0pt] at (1.600,-2.3) {$.$};
\node[anchor=base,inner sep=0pt] at (1.900,-2.3) {$6$};
\node[anchor=base,inner sep=0pt] at (2.700,-2.3) {$4$};
\node[anchor=base,inner sep=0pt] at (-.7,-1.1) {$-$};
\draw[line width=.5pt] (-1,-1.52)--(3.300,-1.52);
\draw[line width=.35pt] (-0.160,.03)--(0.160,.43);
\node[anchor=base,inner sep=0pt,font=\normalsize] at (0.000,.72) {$1$};
\draw[line width=.35pt] (0.640,.03)--(0.960,.43);
\node[anchor=base,inner sep=0pt,font=\normalsize] at (0.800,.72) {$9$};
\draw[line width=.35pt] (1.740,.03)--(2.060,.43);
\node[anchor=base,inner sep=0pt,font=\normalsize] at (1.900,.72) {$9$};
\draw[line width=.35pt] (2.540,.03)--(2.860,.43);
\node[anchor=base,inner sep=0pt,font=\normalsize] at (2.700,.72) {$10$};
\end{tikzpicture}
\par\smallskip{\normalsize 划线数字是退位前的数，上方小字是退位后参与计算的数。\par}
\vspace{0pt}
\begin{enumerate}[leftmargin=1.7em,itemsep=5pt,topsep=0pt]
\item 十位退位：\(\begin{aligned}[t]&\phantom{={}}2-1\\&=1\end{aligned}\)\quad 借出1给后一位。
\item 个位传递退位：\(\begin{aligned}[t]&\phantom{={}}0+10-1\\&=10-1\\&=9\end{aligned}\)\quad 收到10，再借出1给后一位。
\item 十分位传递退位：\(\begin{aligned}[t]&\phantom{={}}0+10-1\\&=10-1\\&=9\end{aligned}\)\quad 收到10，再借出1给后一位。
\item 百分位接收退位：\(\begin{aligned}[t]&\phantom{={}}0+10\\&=10\end{aligned}\)\quad 前一位的1变成本位的10。
\item 百分位：\(\begin{aligned}[t]&\phantom{={}}10-6\\&=4\end{aligned}\)\quad 本位写4。
\item 十分位：\(\begin{aligned}[t]&\phantom{={}}9-3\\&=6\end{aligned}\)\quad 本位写6。
\item 个位：\(\begin{aligned}[t]&\phantom{={}}9-9\\&=0\end{aligned}\)\quad 本位写0。
\item 十位：\(\begin{aligned}[t]&\phantom{={}}1-0\\&=1\end{aligned}\)\quad 本位写1。
\end{enumerate}
\par\smallskip
因此，\(\begin{aligned}[t]&\phantom{={}}20-9.36\\&=10.64\end{aligned}\)。
\par\medskip

\par\medskip\noindent
\textbf{第9小题：\(8.35+4.65=13\)}\par\smallskip
\noindent
\begin{tikzpicture}[x=1em,y=1em,baseline=(current bounding box.center),draw=answerred,text=answerred]
\node[anchor=base,inner sep=0pt] at (0.800,0) {$8$};
\node[anchor=base,inner sep=0pt] at (1.600,0) {$.$};
\node[anchor=base,inner sep=0pt] at (1.900,0) {$3$};
\node[anchor=base,inner sep=0pt] at (2.700,0) {$5$};
\node[anchor=base,inner sep=0pt] at (0.800,-1.1) {$4$};
\node[anchor=base,inner sep=0pt] at (1.600,-1.1) {$.$};
\node[anchor=base,inner sep=0pt] at (1.900,-1.1) {$6$};
\node[anchor=base,inner sep=0pt] at (2.700,-1.1) {$5$};
\node[anchor=base,inner sep=0pt] at (0.000,-2.3) {$1$};
\node[anchor=base,inner sep=0pt] at (0.800,-2.3) {$3$};
\node[anchor=base,inner sep=0pt] at (1.600,-2.3) {$.$};
\node[anchor=base,inner sep=0pt] at (1.900,-2.3) {$0$};
\node[anchor=base,inner sep=0pt] at (2.700,-2.3) {$0$};
\node[anchor=base,inner sep=0pt] at (-.7,-1.1) {$+$};
\draw[line width=.5pt] (-1,-1.52)--(3.300,-1.52);
\node[anchor=base,inner sep=0pt,font=\normalsize] at (1.900,.72) {$1$};
\node[anchor=base,inner sep=0pt,font=\normalsize] at (0.800,.72) {$1$};
\node[anchor=base,inner sep=0pt,font=\normalsize] at (0.000,.72) {$1$};
\end{tikzpicture}
\par\smallskip{\normalsize 上方小字表示从后一位进来的数。\par}
\vspace{0pt}
\begin{enumerate}[leftmargin=1.7em,itemsep=5pt,topsep=0pt]
\item 百分位：\(\begin{aligned}[t]&\phantom{={}}5+5\\&=10\end{aligned}\)\quad 本位写0，向前一位进1。
\item 十分位：\(\begin{aligned}[t]&\phantom{={}}3+6+1\\&=9+1\\&=10\end{aligned}\)\quad 本位写0，向前一位进1。
\item 个位：\(\begin{aligned}[t]&\phantom{={}}8+4+1\\&=12+1\\&=13\end{aligned}\)\quad 本位写3，向前一位进1。
\end{enumerate}
最高位再写进上来的1。\par
\par\smallskip
因此，\(\begin{aligned}[t]&\phantom{={}}8.35+4.65\\&=13.00\end{aligned}\)。
\par\medskip

\par\medskip\noindent
\textbf{第10小题：\(1.19+7.9=9.09\)}\par\smallskip
\noindent
\begin{tikzpicture}[x=1em,y=1em,baseline=(current bounding box.center),draw=answerred,text=answerred]
\node[anchor=base,inner sep=0pt] at (0.000,0) {$1$};
\node[anchor=base,inner sep=0pt] at (0.800,0) {$.$};
\node[anchor=base,inner sep=0pt] at (1.100,0) {$1$};
\node[anchor=base,inner sep=0pt] at (1.900,0) {$9$};
\node[anchor=base,inner sep=0pt] at (0.000,-1.1) {$7$};
\node[anchor=base,inner sep=0pt] at (0.800,-1.1) {$.$};
\node[anchor=base,inner sep=0pt] at (1.100,-1.1) {$9$};
\node[anchor=base,inner sep=0pt] at (1.900,-1.1) {$0$};
\node[anchor=base,inner sep=0pt] at (0.000,-2.3) {$9$};
\node[anchor=base,inner sep=0pt] at (0.800,-2.3) {$.$};
\node[anchor=base,inner sep=0pt] at (1.100,-2.3) {$0$};
\node[anchor=base,inner sep=0pt] at (1.900,-2.3) {$9$};
\node[anchor=base,inner sep=0pt] at (-.7,-1.1) {$+$};
\draw[line width=.5pt] (-1,-1.52)--(2.500,-1.52);
\node[anchor=base,inner sep=0pt,font=\normalsize] at (0.000,.72) {$1$};
\end{tikzpicture}
\par\smallskip{\normalsize 上方小字表示从后一位进来的数。\par}
\vspace{0pt}
\begin{enumerate}[leftmargin=1.7em,itemsep=5pt,topsep=0pt]
\item 百分位：\(\begin{aligned}[t]&\phantom{={}}9+0\\&=9\end{aligned}\)\quad 本位写9。
\item 十分位：\(\begin{aligned}[t]&\phantom{={}}1+9\\&=10\end{aligned}\)\quad 本位写0，向前一位进1。
\item 个位：\(\begin{aligned}[t]&\phantom{={}}1+7+1\\&=8+1\\&=9\end{aligned}\)\quad 本位写9。
\end{enumerate}
\par\smallskip
因此，\(\begin{aligned}[t]&\phantom{={}}1.19+7.9\\&=9.09\end{aligned}\)。
\par\medskip

\par\medskip\noindent
\textbf{第11小题：\(10-3.28=6.72\)}\par\smallskip
\noindent
\begin{tikzpicture}[x=1em,y=1em,baseline=(current bounding box.center),draw=answerred,text=answerred]
\node[anchor=base,inner sep=0pt] at (0.000,0) {$1$};
\node[anchor=base,inner sep=0pt] at (0.800,0) {$0$};
\node[anchor=base,inner sep=0pt] at (1.600,0) {$.$};
\node[anchor=base,inner sep=0pt] at (1.900,0) {$0$};
\node[anchor=base,inner sep=0pt] at (2.700,0) {$0$};
\node[anchor=base,inner sep=0pt] at (0.800,-1.1) {$3$};
\node[anchor=base,inner sep=0pt] at (1.600,-1.1) {$.$};
\node[anchor=base,inner sep=0pt] at (1.900,-1.1) {$2$};
\node[anchor=base,inner sep=0pt] at (2.700,-1.1) {$8$};
\node[anchor=base,inner sep=0pt] at (0.800,-2.3) {$6$};
\node[anchor=base,inner sep=0pt] at (1.600,-2.3) {$.$};
\node[anchor=base,inner sep=0pt] at (1.900,-2.3) {$7$};
\node[anchor=base,inner sep=0pt] at (2.700,-2.3) {$2$};
\node[anchor=base,inner sep=0pt] at (-.7,-1.1) {$-$};
\draw[line width=.5pt] (-1,-1.52)--(3.300,-1.52);
\draw[line width=.35pt] (-0.160,.03)--(0.160,.43);
\node[anchor=base,inner sep=0pt,font=\normalsize] at (0.000,.72) {$0$};
\draw[line width=.35pt] (0.640,.03)--(0.960,.43);
\node[anchor=base,inner sep=0pt,font=\normalsize] at (0.800,.72) {$9$};
\draw[line width=.35pt] (1.740,.03)--(2.060,.43);
\node[anchor=base,inner sep=0pt,font=\normalsize] at (1.900,.72) {$9$};
\draw[line width=.35pt] (2.540,.03)--(2.860,.43);
\node[anchor=base,inner sep=0pt,font=\normalsize] at (2.700,.72) {$10$};
\end{tikzpicture}
\par\smallskip{\normalsize 划线数字是退位前的数，上方小字是退位后参与计算的数。\par}
\vspace{0pt}
\begin{enumerate}[leftmargin=1.7em,itemsep=5pt,topsep=0pt]
\item 十位退位：\(\begin{aligned}[t]&\phantom{={}}1-1\\&=0\end{aligned}\)\quad 借出1给后一位。
\item 个位传递退位：\(\begin{aligned}[t]&\phantom{={}}0+10-1\\&=10-1\\&=9\end{aligned}\)\quad 收到10，再借出1给后一位。
\item 十分位传递退位：\(\begin{aligned}[t]&\phantom{={}}0+10-1\\&=10-1\\&=9\end{aligned}\)\quad 收到10，再借出1给后一位。
\item 百分位接收退位：\(\begin{aligned}[t]&\phantom{={}}0+10\\&=10\end{aligned}\)\quad 前一位的1变成本位的10。
\item 百分位：\(\begin{aligned}[t]&\phantom{={}}10-8\\&=2\end{aligned}\)\quad 本位写2。
\item 十分位：\(\begin{aligned}[t]&\phantom{={}}9-2\\&=7\end{aligned}\)\quad 本位写7。
\item 个位：\(\begin{aligned}[t]&\phantom{={}}9-3\\&=6\end{aligned}\)\quad 本位写6。
\item 十位：\(\begin{aligned}[t]&\phantom{={}}0-0\\&=0\end{aligned}\)\quad 本位写0。
\end{enumerate}
\par\smallskip
因此，\(\begin{aligned}[t]&\phantom{={}}10-3.28\\&=6.72\end{aligned}\)。
\par\medskip

\par\medskip\noindent
\textbf{第12小题：\(15.3-9.56=5.74\)}\par\smallskip
\noindent
\begin{tikzpicture}[x=1em,y=1em,baseline=(current bounding box.center),draw=answerred,text=answerred]
\node[anchor=base,inner sep=0pt] at (0.000,0) {$1$};
\node[anchor=base,inner sep=0pt] at (0.800,0) {$5$};
\node[anchor=base,inner sep=0pt] at (1.600,0) {$.$};
\node[anchor=base,inner sep=0pt] at (1.900,0) {$3$};
\node[anchor=base,inner sep=0pt] at (2.700,0) {$0$};
\node[anchor=base,inner sep=0pt] at (0.800,-1.1) {$9$};
\node[anchor=base,inner sep=0pt] at (1.600,-1.1) {$.$};
\node[anchor=base,inner sep=0pt] at (1.900,-1.1) {$5$};
\node[anchor=base,inner sep=0pt] at (2.700,-1.1) {$6$};
\node[anchor=base,inner sep=0pt] at (0.800,-2.3) {$5$};
\node[anchor=base,inner sep=0pt] at (1.600,-2.3) {$.$};
\node[anchor=base,inner sep=0pt] at (1.900,-2.3) {$7$};
\node[anchor=base,inner sep=0pt] at (2.700,-2.3) {$4$};
\node[anchor=base,inner sep=0pt] at (-.7,-1.1) {$-$};
\draw[line width=.5pt] (-1,-1.52)--(3.300,-1.52);
\draw[line width=.35pt] (-0.160,.03)--(0.160,.43);
\node[anchor=base,inner sep=0pt,font=\normalsize] at (0.000,.72) {$0$};
\draw[line width=.35pt] (0.640,.03)--(0.960,.43);
\node[anchor=base,inner sep=0pt,font=\normalsize] at (0.800,.72) {$14$};
\draw[line width=.35pt] (1.740,.03)--(2.060,.43);
\node[anchor=base,inner sep=0pt,font=\normalsize] at (1.900,.72) {$12$};
\draw[line width=.35pt] (2.540,.03)--(2.860,.43);
\node[anchor=base,inner sep=0pt,font=\normalsize] at (2.700,.72) {$10$};
\end{tikzpicture}
\par\smallskip{\normalsize 划线数字是退位前的数，上方小字是退位后参与计算的数。\par}
\vspace{0pt}
\begin{enumerate}[leftmargin=1.7em,itemsep=5pt,topsep=0pt]
\item 十分位退位：\(\begin{aligned}[t]&\phantom{={}}3-1\\&=2\end{aligned}\)\quad 借出1给后一位。
\item 百分位接收退位：\(\begin{aligned}[t]&\phantom{={}}0+10\\&=10\end{aligned}\)\quad 前一位的1变成本位的10。
\item 百分位：\(\begin{aligned}[t]&\phantom{={}}10-6\\&=4\end{aligned}\)\quad 本位写4。
\item 个位退位：\(\begin{aligned}[t]&\phantom{={}}5-1\\&=4\end{aligned}\)\quad 借出1给后一位。
\item 十分位接收退位：\(\begin{aligned}[t]&\phantom{={}}2+10\\&=12\end{aligned}\)\quad 前一位的1变成本位的10。
\item 十分位：\(\begin{aligned}[t]&\phantom{={}}12-5\\&=7\end{aligned}\)\quad 本位写7。
\item 十位退位：\(\begin{aligned}[t]&\phantom{={}}1-1\\&=0\end{aligned}\)\quad 借出1给后一位。
\item 个位接收退位：\(\begin{aligned}[t]&\phantom{={}}4+10\\&=14\end{aligned}\)\quad 前一位的1变成本位的10。
\item 个位：\(\begin{aligned}[t]&\phantom{={}}14-9\\&=5\end{aligned}\)\quad 本位写5。
\item 十位：\(\begin{aligned}[t]&\phantom{={}}0-0\\&=0\end{aligned}\)\quad 本位写0。
\end{enumerate}
\par\smallskip
因此，\(\begin{aligned}[t]&\phantom{={}}15.3-9.56\\&=5.74\end{aligned}\)。
\par\medskip

}

\QHeading{第3题　解决问题。}

\Problem{① 小明身高1.45米，小红比小明矮0.08米，小红身高多少米？}
\par


\MethodTitle{方法一}
\Answer{$1.45-0.08=1.37\text{（米）}$\par 答：小红身高1.37米。}

\MethodTitle{方法二：多减再补回}
\Answer{把0.08米先按0.1米减，多减了0.02米。}
\Answer{\Formula{\begin{aligned}[t]&\phantom{={}}1.45-0.08\\
&=1.45-0.1+0.02\\
&=1.35+0.02\\
&=1.37\text{（米）}\end{aligned}}}
\Answer{答：小红身高1.37米。}
\Problem{② 一桶油重5.5千克，用去2.85千克，还剩多少千克？}

\MethodTitle{方法一}
\Answer{$5.5-2.85=2.65\text{（千克）}$\par 答：还剩2.65千克。}

\MethodTitle{方法二：把用油量凑成3千克}
\Answer{先减3千克，多减0.15千克，再补回。}
\Answer{\Formula{\begin{aligned}[t]&\phantom{={}}5.5-2.85\\
&=5.5-3+0.15\\
&=2.5+0.15\\
&=2.65\text{（千克）}\end{aligned}}}
\Answer{答：还剩2.65千克。}

